\documentclass[11pt,a4paper]{article}
\usepackage[utf8]{inputenc}
\usepackage[T1]{fontenc}
\usepackage[margin=0.75in]{geometry}
\usepackage{enumitem}
\usepackage{titlesec}
\usepackage{xcolor}
\usepackage{hyperref}
\definecolor{primarycolor}{RGB}{33,37,41}
\definecolor{secondarycolor}{RGB}{0,102,204}
\definecolor{accentcolor}{RGB}{108,117,125}
\hypersetup{colorlinks=true,linkcolor=secondarycolor,urlcolor=secondarycolor,citecolor=secondarycolor}
\titleformat{\section}{\large\bfseries\color{primarycolor}}{}{0em}{}
\titlespacing*{\section}{0pt}{0.8em}{0.35em}
\newcommand{\cvsection}[1]{\section*{#1}}
\setlist[itemize]{leftmargin=*,itemsep=0.18em,topsep=0.25em,parsep=0pt}
\setlength{\parindent}{0pt}
\setlength{\parskip}{0.25em}
\pagestyle{empty}
\begin{document}
{\LARGE\bfseries\color{primarycolor} Katarzyna Wiśniewska}\par
{\large\color{accentcolor} Pflegehelferin}\par
\cvsection{PROFIL}
Pflegehelferin mit Erfahrung in der stationären Altenpflege und fünfjähriger Tätigkeit als Krankenschwester auf einer Station für Innere Medizin in Polen. Erfahrung mit Menschen mit Demenz, Sturzprophylaxe und der Einarbeitung neuer Kolleginnen. LG1-Kurs, Deutschzertifikat B2 und gute Kenntnisse in Vivendi PD.
\cvsection{BERUFSERFAHRUNG}
\textbf{Seit 2021 | Pflegehelferin}\par
Seniorenzentrum Am Lindenhof
\begin{itemize}
\item Seit 2022: Leitung einer wöchentlichen Singgruppe für 6 bis 8 Menschen mit Demenz mit polnischen und deutschen Volksliedern; positive Rückmeldung von Angehörigen im Heimbeirat.
\item Seit 2021: Einarbeitung von bisher 7 neuen Pflegehelferinnen und Praktikanten während ihrer ersten zwei Wochen nach dem Einarbeitungsplan der Einrichtung.
\item 2023: Gemeinsam mit der Wohnbereichsleitung Sturzprotokolle aus sechs Monaten ausgewertet und Nachtlichter sowie feste Toilettengänge für 9 Bewohner eingeführt. Im folgenden Halbjahr sanken die Stürze im Wohnbereich von 14 auf 8.
\end{itemize}
\textbf{Fünfjährige Tätigkeit | Krankenschwester, Innere Medizin}\par
Szpital Opole, Polen
\begin{itemize}
\item 2016–2019: Verbandswechsel bei Dekubitus und diabetischem Fuß unter Aufsicht der Stationsärztin.
\end{itemize}
\cvsection{QUALIFIKATIONEN UND ANERKENNUNG}
\begin{itemize}
\item Polnischer Pflegeabschluss im deutschen Anerkennungsverfahren: Defizitbescheid liegt vor; Kenntnisprüfung für März 2027 vorgesehen.
\item Aktueller beruflicher Status in Deutschland: Pflegehelferin, keine Anerkennung als Pflegefachfrau.
\item Kurs Behandlungspflege Leistungsgruppe 1 (LG1).
\end{itemize}
\cvsection{KENNTNISSE UND STÄRKEN}
\begin{itemize}
\item Stationäre Altenpflege und Umgang mit Menschen mit Demenz.
\item Pflegedokumentation: gute Kenntnisse in Vivendi PD.
\item Geduld im Umgang mit älteren Menschen und Freude an Teamarbeit.
\item Deutsch: B2-Zertifikat, 2021.
\end{itemize}
\cvsection{ARBEITSZEIT UND MOBILITÄT}
\begin{itemize}
\item Gewünschter Beschäftigungsumfang: Teilzeit, bevorzugt 75 \%.
\item Einsatz an jedem zweiten Wochenende möglich.
\item Führerschein Klasse B noch nicht vorhanden; Theorieprüfung bestanden, praktische Prüfung für Januar 2027 vorgesehen.
\item Derzeit mobil mit Fahrrad oder E-Bike.
\end{itemize}
\end{document}